%% Distinct first-order calculi additions retained after exact source restoration.
\noindent\textit{Reader note.} The following equality, Hilbert, tableau, and
comparison material is separately attributed editorial apparatus.  It does
not alter, normalize, or replace the exact inherited natural-deduction or
sequent-rule displays above.  The equality laws have the supplied annex
derivation; the Hilbert material is attributed to
\cite{Mendelson1997,Kleene1952,Enderton2001}, with Mendelson supplying the
universal base and Kleene the primitive-existential schema/rule; tableau
adequacy is attributed to
\cite{Smullyan1968}.  The separately identified registry sheets provide
navigation and variant context, not source substitutes; their added policies
and notation do not change the exact inherited calculi.

\section{Equality metatheory}
\label{sec:p02-fol-equality-metatheory}

\begin{hbtheorem}{thm:p02-equality-laws}{proved=annex:p2.fol.calc.4}{Derived equality laws}
For all \(\mathcal L\)-terms \(r,t,u\), every formula \(\varphi\), and every
variable \(x\), the natural-deduction equality rules on
\hyperref[sheet:nd-fol]{\texttt{sheet:nd-fol}} derive:
\[
 \varnothing\derives_{\mathsf{ND}_{\mathsf{FOL}}}t=t,
 \qquad
 \{r=t\}\derives_{\mathsf{ND}_{\mathsf{FOL}}}t=r,
\]
\[
 \{r=t,t=u\}\derives_{\mathsf{ND}_{\mathsf{FOL}}}r=u,
 \qquad
 \{r=t,\varphi[r/x]\}
   \derives_{\mathsf{ND}_{\mathsf{FOL}}}\varphi[t/x].
\]
Every displayed substitution is the total capture-avoiding operation of the
syntax chapter.  Reflexivity and replacement are the registered primitive
rules; symmetry and transitivity are derived, not silently added axioms.
\end{hbtheorem}

The free-standing term-congruence result used by the quotient term model is
Lemma~\ref{lem:equality-congruence}.  Together with the preceding
derived laws, it supplies the consolidated equality metatheory used by the
following proof systems.  The symmetric replacement instance for
\hyperref[sheet:lk-eq]{\texttt{sheet:lk-eq}} is likewise derived by symmetry
and Cut; the sheet comments record that it is not primitive.

\section{An equality-free Hilbert calculus}
\label{sec:p02-hilbert-fol}

\begin{hbdefinition}{def:p02-hilbert-fol-dependencies}{Hilbert derivation with dependencies}
Let \(\mathsf H_{\mathsf{FOL}}\) be
\hyperref[sheet:hilbert-fol]{\texttt{sheet:hilbert-fol}}, whose language is
the equality-free \(\{\neg,\to,\forall,\exists\}\)-fragment.  A
\emph{dependency-annotated derivation from \(\Gamma\)}
\index{Hilbert calculus!first-order}
\index{derivation!dependency-annotated} is a finite sequence
of pairs \((\chi_i,D_i)\), where \(D_i\subseteq\Gamma\) is finite, with
exactly these clauses:
\begin{enumerate}
  \item if \(\chi_i\in\Gamma\), then \(D_i=\{\chi_i\}\);
  \item an axiom-schema instance has \(D_i=\varnothing\);
  \item modus ponens from lines \(j,k<i\) has
        \(D_i=D_j\cup D_k\); and
  \item generalization from \((\chi_j,D_j)\) to
        \((\forall x\,\chi_j,D_j)\) is permitted exactly when
        \(x\notin\FV(\delta)\) for every \(\delta\in D_j\); and
  \item existential elimination from
        \((\theta\to\psi,D_j)\) to
        \((\exists x\,\theta\to\psi,D_j)\) is permitted exactly when
        \(x\notin\FV(\psi)\) and
        \(x\notin\FV(\delta)\) for every \(\delta\in D_j\).
\end{enumerate}
Write \(\Gamma\derives_{\mathsf H_{\mathsf{FOL}}}\chi\) iff some such
sequence ends in \((\chi,D)\).  The annotation is metalinguistic and
eliminable: erasing it gives the sheet's finite formula sequence, while
reconstructing it records exactly which premises license each generalization.
It is needed below because the first-order deduction theorem fails for an
unrestricted generalization rule.
\end{hbdefinition}

\begin{hbtheorem}{thm:p02-hilbert-fol-deduction}{proved=annex:p2.fol.calc.5}{Hilbert-FOL deduction theorem}
For every set \(\Gamma\) and formulas \(\varphi,\psi\) of the sheet's
fragment,
\[
 \Gamma\cup\{\varphi\}
   \derives_{\mathsf H_{\mathsf{FOL}}}\psi
 \quad\Longleftrightarrow\quad
 \Gamma\derives_{\mathsf H_{\mathsf{FOL}}}(\varphi\to\psi).
\]
The derivability relation is the dependency-annotated one just defined; the
theorem does not license generalizing a variable free in a premise on which
the generalized line depends.
\end{hbtheorem}

Mendelson supplies the universal Hilbert base and dependency-sensitive
generalization convention; Kleene supplies existential-introduction schema 11
and existential-elimination rule 12, in the premise-sensitive form displayed
on the sheet; see \cite{Mendelson1997,Kleene1952,Enderton2001}.  The
equality-free fragment must not be
reported as literally the same language as the equality-bearing natural-
deduction and sequent sheets.

\section{First-order analytic tableaux}
\label{sec:p02-tableaux-fol}

\begin{hbdefinition}{def:p02-fol-tableau}{First-order tableau}
For a finite set \(\Gamma\) of equality-free \(\mathcal L\)-sentences, let
\(\mathcal L_{\Gamma}\) contain exactly the nonlogical symbols occurring in
\(\Gamma\), and obtain \(\mathcal L_{\Gamma}^*\) by adjoining new parameter
constants \(a_n\) for \(n\in\mathbb N\).  The resulting signature has only
countably many closed terms; fix an enumeration of all of them.  A
\emph{first-order tableau rooted at \(\Gamma\)}\index{tableau!first-order}
is a finitely branching
retained-formula tree generated by
\hyperref[sheet:tableaux-fol]{\texttt{sheet:tableaux-fol}}.  Propositional
and witness-producing occurrences expand once.  A universal occurrence (or
the negation of an existential) records each closed
\(\mathcal L_{\Gamma}^*\)-term already used and remains available at every
unused term.  A branch is \emph{fair} iff every such occurrence is eventually
instantiated with every term in the fixed enumeration, whether or not that
term previously appeared on the branch.  A branch is \emph{closed} iff it
contains \(\bottom\) or
both \(\varphi\) and \(\neg\varphi\) for some \(\varphi\), and is \emph{open}
otherwise; it is \emph{saturated} iff no expansion applicable on that branch
is omitted.  These occurrence marks and parameters are
metalinguistic proof-search data, not new symbols of the original language.
\end{hbdefinition}

\begin{hbtheorem}{thm:p02-fol-tableau-adequacy}{cited=Smullyan1968}{First-order tableau adequacy and countermodel extraction}
Let \(\Gamma\) be a finite set of equality-free \(\mathcal L\)-sentences.
Then \(\Gamma\) is unsatisfiable iff a finite closed tableau rooted at
\(\Gamma\) exists.  Moreover, for every fair saturated open branch \(B\),
define its positive atomic part by
\[
 B^{+}_{\mathrm{at}}
 \defeq
 \{R(t_1,\ldots,t_n):R(t_1,\ldots,t_n)\text{ occurs on }B\}.
\]
The corresponding Herbrand structure is
\(\mathcal M_{B^{+}_{\mathrm{at}}}\).  Its domain is the closed
\(\mathcal L_{\Gamma}^*\)-terms, its functions act by term formation, and
its relations are exactly the members of \(B^{+}_{\mathrm{at}}\).
The branch truth induction gives
\[
 \mathcal M_{B^{+}_{\mathrm{at}}}\entails B.
\]
Hence its
\(\mathcal L_{\Gamma}\)-reduct models \(\Gamma\), and interpreting every
unused \(\mathcal L\)-symbol arbitrarily expands it to an
\(\mathcal L\)-model of \(\Gamma\).
This is the countermodel-extraction clause, not merely a closure test.
\end{hbtheorem}

The equality-free scope is load-bearing.  Identity tableaux require an
additional equality calculus; neither the equality rules of
\hyperref[sheet:nd-fol]{\texttt{sheet:nd-fol}} nor those of
\hyperref[sheet:lk-eq]{\texttt{sheet:lk-eq}} are silently available
on the tableau sheet.  Together with the already verified propositional entry,
this supplies soundness, completeness, and open-branch countermodel extraction
for both propositional and equality-free first-order tableaux.

\begin{landscape}
\section{First-order proof-system comparison}
\label{sec:p02-fol-calculi-comparison}

\begingroup\scriptsize\setlength{\tabcolsep}{3pt}
\begin{longtable}{@{}p{.65in}p{1.35in}p{1.95in}p{1.90in}p{1.25in}@{}}
\caption{First-order proof systems side by side.}
\label{tab:p02-fol-proof-systems}\\
\toprule
System & Rule families & Exact side conditions & Adequacy evidence & Scope and registry home \\ \midrule
\endfirsthead
\toprule
System & Rule families & Exact side conditions & Adequacy evidence & Scope and registry home \\ \midrule
\endhead
Hilbert & Propositional schemas A1--A3; quantifier schemas
\(\forall1,\forall2,\exists1\); rules MP, Gen, and \(\exists E\). & All substitutions are
total and capture-avoiding; in \(\forall2\), \(x\notin\FV(\varphi)\); in
\(\exists E\), \(x\notin\FV(\psi)\) and \(x\) is not free in any premise on
which its premise line depends; Gen has the analogous dependency condition. &
Soundness and strong completeness are obtained from the split sources:
Mendelson supplies A1--A5 and Gen, while Kleene supplies
existential-introduction schema 11 and existential-elimination rule 12.  These
make
\(\exists x\varphi\) interderivable with
\(\neg\forall x\neg\varphi\), so the usual elimination of primitive
\(\exists\) transfers Mendelson's strong-completeness theorem to this combined
surface \cite{Mendelson1997,Kleene1952}; the deduction theorem is proved here as
Theorem~\ref{thm:p02-hilbert-fol-deduction}. & Equality-free
\(\neg,\to,\forall,\exists\) fragment;
\hyperref[sheet:hilbert-fol]{\texttt{sheet:hilbert-fol}}. \\
Natural deduction & Propositional introduction/elimination rules, RAA,
\(\forall I/E\), \(\exists I/E\), alpha-conversion, and \(=I/E\). &
\(\forall I\): \(x\notin\FV(\Gamma)\); \(\exists E\): the witness variable is
absent from the conclusion, all other open assumptions, and the existential
major premise; every discharge is nonvacuous; substitution is total and
capture-avoiding. & Proved soundness,
Theorem~\ref{thm:fol-soundness}, and proved strong completeness,
Theorem~\ref{thm:fol-completeness} and
Corollary~\ref{cor:fol-completeness-countable}. & Primitive identity;
\hyperref[sheet:nd-fol]{\texttt{sheet:nd-fol}}. \\
\(\mathsf{LK}\) & Contexted initials; structural and propositional left/right
rules; Cut and Mix; \(\forall L/R\), \(\exists L/R\); the equality extension
adds reflexivity and left replacement. & \(\forall R\) and \(\exists L\) use
an eigenvariable free in no context or principal formula; \(\forall L\) and
\(\exists R\) use total capture-avoiding substitution; Mix has positive
multiplicities. & Proved soundness,
Theorem~\ref{thm:fol-lk-soundness}; honest completeness sketch,
Theorem~\ref{thm:fol-lk-completeness}. & Equality-free
\hyperref[sheet:lk]{\texttt{sheet:lk}}; identity extension
\hyperref[sheet:lk-eq]{\texttt{sheet:lk-eq}}. \\
Tableau & Propositional linear/branching expansions; repeated gamma and
one-use delta expansions; complementary-pair or \(\bottom\) closure. & Finite
equality-free sentence root; gamma fairness ranges over every enumerated
closed term; delta uses the least parameter not earlier on the branch; every
substitution is total and capture-avoiding. & Cited finite closed-refutation
completeness and fair-open-branch model extraction,
Theorem~\ref{thm:p02-fol-tableau-adequacy}. & Equality-free;
\hyperref[sheet:tableaux-fol]{\texttt{sheet:tableaux-fol}}. \\
\bottomrule
\end{longtable}
\endgroup
\end{landscape}
